\section*{الفصل \ref{ch:g12:glucose-and-diabetes} --- سكر الدم وداء السكري}

\begin{solution}{exo:g12:glucose-and-diabetes:1}
نحو \qty{1}{g/L}، أي \qty{5.5}{mmol/L}. ودون \qty{0.5}{g/L} يخفق
الدماغ؛ وفوق \qty{1.8}{g/L} يمر الغلوكوز إلى البول (وسنوات فوق
\qty{1.3}{g/L} تتلف الأوعية).
\end{solution}

\begin{solution}{exo:g12:glucose-and-diabetes:2}
\omterm{prop:g12:glucose-and-diabetes:hormones}{الأنسولين}، من \omterm{def:g10:cells-common-unit:cell}{الخلايا} $\beta$، يخفض \omterm{def:g12:glucose-and-diabetes:glycaemia}{سكر الدم} (بأخذه وادخاره في العضلة
والدهن والكبد)؛ \omterm{prop:g12:glucose-and-diabetes:hormones}{والغلوكاغون}، من \omterm{def:g10:cells-common-unit:cell}{الخلايا} $\alpha$، يرفعه (بتحرير الكبد
الغلوكوز).
\end{solution}

\begin{solution}{exo:g12:glucose-and-diabetes:3}
الكبد، ويدخره غليكوجينًا (وهو يستطيع أيضًا صنع غلوكوز جديد من \omterm{def:g10:chemistry-of-life:families}{الحموض
الأمينية} واللاكتات).
\end{solution}

\begin{solution}{exo:g12:glucose-and-diabetes:4}
النمط 1: \omterm{def:g10:cells-common-unit:cell}{الخلايا} $\beta$ مخربة \omterm{prop:g12:glucose-and-diabetes:hormones}{والأنسولين} غائب. والنمط 2: \omterm{prop:g12:glucose-and-diabetes:hormones}{الأنسولين}
موجود لكن الأنسجة تستجيب له استجابة ضعيفة، وتتعب \omterm{def:g10:cells-common-unit:cell}{الخلايا} $\beta$ في
النهاية.
\end{solution}

\begin{solution}{exo:g12:glucose-and-diabetes:5}
السليم نحو \qty{0.95}{g/L}؛ والتحمل المضطرب نحو \qty{1.7}{g/L}
(بين 1.4 و2)؛ والنمط 2 نحو \qty{2.7}{g/L} (فوق 2).
\end{solution}

\begin{solution}{exo:g12:glucose-and-diabetes:6}
الكبد يحرر الغلوكوز من غليكوجينه بالمعدل الذي ينزعه به الدماغ، بضبط
\omterm{prop:g12:glucose-and-diabetes:hormones}{الغلوكاغون}، فيبقى المخزن قرب \qty{5}{g} مع أن محتواه يتجدد كل ساعة.
\end{solution}

\begin{solution}{exo:g12:glucose-and-diabetes:7}
القناة تحمل \omterm{def:g11:enzymes-and-phenotype:enzyme}{الإنزيمات} الهاضمة إلى المعي؛ وسدها يترك \omterm{def:g12:glucose-and-diabetes:glycaemia}{سكر الدم} طبيعيًا،
فأثر البنكرياس في الغلوكوز ليس عبر الهضم. أما نزع العضو كله فينزع
\omterm{prop:g12:glucose-and-diabetes:hormones}{الجزر}: \omterm{prop:g12:glucose-and-diabetes:hormones}{فالأنسولين} يصنع في البنكرياس ويبلغ أهدافه عبر الدم لا عبر
القناة.
\end{solution}

\begin{solution}{exo:g12:glucose-and-diabetes:8}
\omterm{prop:g12:glucose-and-diabetes:hormones}{الغلوكاغون} يعمل بأن يجعل الكبد يفكك غليكوجينه؛ وبعد ليلة من الصيام
يكون المخزون حاضرًا فيرتفع \omterm{def:g12:glucose-and-diabetes:glycaemia}{سكر الدم}؛ وبعد ثلاثة أيام يكون \omterm{def:g10:exercise-and-energy:glycogen}{الغليكوجين}
قد استنفد ولا يجد \omterm{prop:g12:glucose-and-diabetes:hormones}{الغلوكاغون} شيئًا يحرره.
\end{solution}

\begin{solution}{exo:g12:glucose-and-diabetes:9}
فوق نحو \qty{1.8}{g/L} لا تستطيع الكليتان أن تعيدا امتصاص كل الغلوكوز
الذي ترشحانه؛ فيمر إلى البول ويجر معه الماء بالتناضح. والماء المفقود
لا بد أن يعوَّض: فالعطش.
\end{solution}

\begin{solution}{exo:g12:glucose-and-diabetes:10}
\omterm{prop:g12:glucose-and-diabetes:hormones}{الأنسولين} يدفع الغلوكوز إلى \omterm{def:g10:cells-common-unit:cell}{خلايا} العضلة والدهن ويوقف تحرير الكبد،
لكن لا يدخل غلوكوز: فيهبط \omterm{def:g12:glucose-and-diabetes:glycaemia}{سكر الدم} دون \qty{0.6}{g/L} في غضون ساعة أو
ساعتين، مع تعرق ورجفان وتخليط؛ ويصححه سكر يؤخذ في الحال.
\end{solution}

\begin{solution}{exo:g12:glucose-and-diabetes:11}
العضلة العاملة تأخذ الغلوكوز من الدم بطريق لا يحتاج \omterm{prop:g12:glucose-and-diabetes:hormones}{الأنسولين}،
والعضلة المدرَّبة تستعيد حساسيتها للأنسولين بين الحصص؛ فيهبط \omterm{def:g12:glucose-and-diabetes:glycaemia}{سكر الدم}
أثناء الرياضة وبعدها، وتقل مقاومة الأنسجة، وهي سبب المرض.
\end{solution}

\begin{solution}{exo:g12:glucose-and-diabetes:12}
لا أنسولين: النمط 1، \omterm{def:g10:cells-common-unit:cell}{فالخلايا} $\beta$ في \omterm{prop:g12:glucose-and-diabetes:hormones}{الجزر} مخربة ولا تفرز شيئًا.
وثلاثة أمثال \omterm{prop:g12:glucose-and-diabetes:hormones}{الأنسولين} الطبيعي: النمط 2، \omterm{prop:g12:glucose-and-diabetes:hormones}{فالجزر} تفرز بأقصى طاقتها في
وجه أنسجة لم تعد تستجيب.
\end{solution}

\begin{solution}{exo:g12:glucose-and-diabetes:13}
السليم: يرتفع \omterm{def:g12:glucose-and-diabetes:glycaemia}{سكر الدم} إلى نحو \qty{1.3}{g/L} عند 30 إلى 60 دقيقة
ويعود في ساعتين؛ ويرتفع \omterm{prop:g12:glucose-and-diabetes:hormones}{الأنسولين} معه ويهبط معه. والنمط 1 بلا علاج:
يتسلق \omterm{def:g12:glucose-and-diabetes:glycaemia}{سكر الدم} فوق \qty{2}{g/L} ويبقى؛ \omterm{prop:g12:glucose-and-diabetes:hormones}{والأنسولين} مستو عند الصفر ---
فلا شيء ينزع الغلوكوز. والنمط 2: يرتفع \omterm{def:g12:glucose-and-diabetes:glycaemia}{سكر الدم} ارتفاعًا كبيرًا ويهبط
ببطء؛ ويرتفع \omterm{prop:g12:glucose-and-diabetes:hormones}{الأنسولين} فوق الطبيعي ويبقى مرتفعًا --- \omterm{def:g11:hormones-and-reproduction:hormone}{هرمون} كثير وأثر
قليل.
\end{solution}

\begin{solution}{exo:g12:glucose-and-diabetes:14}
$5 + 80 = \qty{85}{g}$ في \qty{5}{L}: أي \qty{17}{g/L}. ولا يحدث ذلك
قط لأن الامتصاص يأخذ ساعة أو ساعتين، ولأن \omterm{prop:g12:glucose-and-diabetes:hormones}{الأنسولين}، كلما وصل غلوكوز،
أرسله إلى \omterm{def:g10:exercise-and-energy:glycogen}{غليكوجين} الكبد \omterm{def:g10:exercise-and-energy:glycogen}{وغليكوجين} العضلة والدهن، بينما يحرق الدماغ
والأنسجة بعضه: فلا يحمل المخزن أكثر من بضعة غرامات فوق الطبيعي.
\end{solution}

\begin{solution}{exo:g12:glucose-and-diabetes:15}
كلاهما: مستقبل (\omterm{def:g10:cells-common-unit:cell}{خلايا} \omterm{prop:g12:glucose-and-diabetes:hormones}{الجزر}؛ وتحت المهاد والنخامى)، ونقطة ضبط،
ومنفذات، \omterm{def:g11:hormones-and-reproduction:hormone}{وتغذية راجعة} سلبية تزيل الإشارة. ولسكر الدم هرمونان متعاكسان
يعملان على الأعضاء نفسها، فيستطاع دفع المتغير إلى أعلى وإلى أسفل
دفعًا فاعلًا؛ أما حلقة التستوستيرون فلها \omterm{def:g11:hormones-and-reproduction:hormone}{هرمون} واحد وتصحح بالمزيد منه
أو بالأقل غالبًا. والمنفذان المتعاكسان يعطيان ضبطًا أسرع وأضيق لمتغير
يتغير سريعًا في الاتجاهين.
\end{solution}

\begin{solution}{pb:g12:glucose-and-diabetes:1}
\textbf{1.} \qty{5}{g}؛ و$\qty{1}{g/L} \div \qty{180}{g/mol} = \qty{5.55}{mmol/L}$.

\textbf{2.} ساعة واحدة.

\textbf{3.} $100/5 = 20$ ساعة. وبعدها يصنع الكبد غلوكوزًا جديدًا من
\omterm{def:g10:chemistry-of-life:families}{الحموض الأمينية} (ومن اللاكتات)، على حساب بروتينات الجسم.

\textbf{4.} \omterm{def:g10:cells-common-unit:cell}{خلايا} العضلة تفتقر إلى \omterm{def:g11:enzymes-and-phenotype:enzyme}{الإنزيم} الذي يحرر الغلوكوز الحر من
\omterm{def:g10:exercise-and-energy:glycogen}{الغليكوجين} إلى الدم؛ فغليكوجينها يحرق في مكانه، لتقلصها هي.

\textbf{5.} $5 + 60 = \qty{65}{g}$ في \qty{5}{L}: أي \qty{13}{g/L}.
\omterm{prop:g12:glucose-and-diabetes:hormones}{والأنسولين} يرسل الغلوكوز إلى \omterm{def:g10:exercise-and-energy:glycogen}{غليكوجين} الكبد والعضلة وإلى الدهن، وتحرق
الأنسجة بعضه؛ فيبلغ \omterm{def:g12:glucose-and-diabetes:glycaemia}{سكر الدم} ذروته نحو \qty{1.4}{g/L}.

\textbf{6.} صائمًا \qty{1.15}{g/L} (دون 1.26) لكن \qty{1.85}{g/L} عند
120 دقيقة (بين 1.4 و2): أي تحمل مضطرب، على عتبة \omterm{def:g12:glucose-and-diabetes:diabetes}{داء السكري}.

\textbf{7.} لا: \omterm{prop:g12:glucose-and-diabetes:hormones}{فالأنسولين} عند ضعف ذروة سليمة. وإنما استجابة الأنسجة
له هي المخفقة --- أي مقاومة \omterm{prop:g12:glucose-and-diabetes:hormones}{الأنسولين}.

\textbf{8.} النمط 2 (أو طور إنذاره). والمتوقع: بالغ، زائد الوزن، قليل
الحركة، وكثيرًا ما يكون له أقارب سكريون.

\textbf{9.} سكر دم يبدأ فوق \qty{1.5}{g/L}، ويرتفع فوق \num{2.5} ولا
يهبط؛ وأنسولين قرب الصفر طوال الوقت.

\textbf{10.} إنقاص الوزن والرياضة المنتظمة: فهما يعيدان حساسية الأنسجة
للأنسولين ويتيحان للعضلة العاملة أن تأخذ الغلوكوز من دونه.

\textbf{11.} $90/10 = 9$ وحدات.

\textbf{12.} الجرعة تتصرف بكمية \qty{90}{g} لكن \qty{45}{g} لا غير تصل:
\omterm{prop:g12:glucose-and-diabetes:hormones}{فالأنسولين} الفائض يهبط \omterm{def:g12:glucose-and-diabetes:glycaemia}{بسكر الدم} بما يعادل \qty{45}{g} في المخزن
والمدخرات --- أي دون \qty{0.6}{g/L} بكثير: تعرق، ورجفان، وتخليط؛
وسكر في الحال.

\textbf{13.} $(5 + 90)/5 = \qty{19}{g/L}$ لو لم ينزع شيء؛ وفي الواقع
عدة \unit{g/L}. وفوق \qty{1.8}{g/L} تدع الكليتان الغلوكوز يمر إلى
البول، ومعه الماء: بول غزير وعطش.

\textbf{14.} بين الوجبات كان الكبد سيحرر الغلوكوز بلا كابح (\omterm{prop:g12:glucose-and-diabetes:hormones}{فالغلوكاغون}
يعمل ولا شيء يعارضه) وكانت الأنسجة ستحرق الدسم بنواتجه الحمضية: فخلفية
من \omterm{prop:g12:glucose-and-diabetes:hormones}{الأنسولين} لازمة للتوازن الأساسي، لا للوجبات وحدها.

\textbf{15.} المقياس يقيس \omterm{def:g12:glucose-and-diabetes:glycaemia}{سكر الدم} كما كانت تقيسه \omterm{def:g10:cells-common-unit:cell}{الخلية} $\beta$؛
والقلم يسلّم \omterm{prop:g12:glucose-and-diabetes:hormones}{الأنسولين} الذي كانت ستفرزه. وهما لا يستطيعان تعويض
التعديل المتصل دقيقةً دقيقةً: فالمريضة تقيس بضع مرات في اليوم وتقدّر،
حيث كانت \omterm{def:g10:cells-common-unit:cell}{الخلية} تقيس باستمرار وتستجيب بالضبط.

\textbf{16.} النمط 2، فخلاياه $\beta$ موجودة ويمكن دفعها. أما في النمط
1 فلا \omterm{def:g10:cells-common-unit:cell}{خلايا} $\beta$ تنبَّه.

\textbf{17.} استجابة الأنسجة للأنسولين --- أي جانب المنفذات من الحلقة:
فمع دهن أقل وعضلة نشطة، تعود إشارة أنسولين عادية تنتج أخذًا عاديًا،
وترتاح \omterm{def:g10:cells-common-unit:cell}{الخلايا} $\beta$.

\textbf{18.} كل منهما يفرز استجابةً للانحراف المعاكس في المتغير نفسه:
فارتفاع \omterm{def:g12:glucose-and-diabetes:glycaemia}{سكر الدم} يستدعي \omterm{prop:g12:glucose-and-diabetes:hormones}{الأنسولين} ويسكت \omterm{prop:g12:glucose-and-diabetes:hormones}{الغلوكاغون}، وانخفاضه بالعكس.
وعند \qty{1}{g/L} تفرز الجزيرة كليهما بمعدل منخفض متوازن --- وهي حالة
راحتها.

\textbf{19.} انخفاض \omterm{def:g12:glucose-and-diabetes:glycaemia}{سكر الدم} يقتل الدماغ في غضون دقائق، وارتفاعه يتلف
ببطء عبر سنين: فللخطر الحاد عدة حراس مستقلين، وللبطيء حارس واحد ---
ولهذا كان إخفاق \omterm{prop:g12:glucose-and-diabetes:hormones}{الأنسولين} وحده كافيًا لإحداث مرض، ولهذا كان \omterm{def:g12:glucose-and-diabetes:diabetes}{داء
السكري} شائعًا.

\textbf{20.} نحو \qty{5}{g} من الغلوكوز في الدم؛ \omterm{def:g10:exercise-and-energy:glycogen}{وغليكوجين} الكبد يغطي
نحو 20 ساعة؛ ومريض القسم الثاني عنده تحمل مضطرب على حافة النمط 2 من
\omterm{def:g12:glucose-and-diabetes:diabetes}{داء السكري}؛ وجرعة القسم الثالث 9 وحدات.
\end{solution}
